\documentclass[11pt]{article}
\usepackage[a4paper,margin=25mm]{geometry}
\usepackage{amsmath,amssymb,amsthm,booktabs,hyperref}
\newtheorem{proposition}{Proposition}
\newtheorem{lemma}{Lemma}
\newcommand{\Z}{\mathbb Z}
\title{Construction searches for distinct covering systems\\General and squarefree moduli}
\author{Research working note}
\date{15 October 2026}
\begin{document}
\maketitle

\paragraph{Status.}
This investigation has \emph{not} produced a covering system with minimum
modulus greater than 42. It supplies a finite search formulation, a proof
that successful certificates yield finite distinct covering systems, an exact
checker, and bounded initial experiments. The positive controls have minimum
moduli 3, 6, and 10; a new verified finite-completion schema reaches 20.
They test the machinery and are not record improvements.
The thesis-based continuation gives a restricted spare-family obstruction
and six finite partial-cover schemas covering $3/10$ of the whole
$10\pmod{42}$ branch, but does not complete its replacement.
The user subsequently raised the construction target to minimum modulus
at least 1000. That target, and even the target 43, remain unresolved here.
The target-100 continuation uses exact overlap counts and
fresh-prime divisor blocks. It produces an independently rejected candidate
and a rigorous obstruction for that restricted repair family, not a new cover.
At target 60, five exact candidates fail independent verification;
a new policy forbidding overlap between kernel cylinders also fails to
complete a cover. These experiments do not prove that minimum 60 is impossible.
The squarefree target-30 experiment used a separate ordinary-residue model
and failed; a squarefree minimum-2 control is verified. At squarefree minimum
3, a joint parity construction and exact pair repairs
give verified partial systems, not a complete cover. A normalization argument
also excludes all prime supports contained in the primes through 11.
The current active target is unrestricted minimum modulus 50. Exact nested
resource exchanges and graded arrow blocks have been implemented and tested.
No minimum-50 cover has been found. The branch-dependent extension now has
verified partial replacement schemas, with explicitly checked dependencies.

\paragraph{Sources and scope.}
Owens's 2014 thesis establishes minimum modulus 42~\cite{owens}.
The full thesis was obtained on 15 October 2026 through the institutional
page's browser download, after direct PDF requests returned HTTP 403.
Sections 3.1--3.4 and the later dependency structure have now been examined.
The whole 42 construction has not been independently machine-verified.
Nielsen's 2009 paper~\cite{nielsen} was read
from the user's local literature collection, including Sections 3--5.
It describes recursive prime-power tails and their finite closure, and explains
why removing the smallest class is not repaired simply by introducing later
primes. The following formulation uses that finite-tail idea. No publication
priority or claim of a new record is made.

\section{A finite model retaining prime-power tails}
Fix distinct primes $p_1,\ldots,p_s$ and integers $E_i\geq1$.
For the coordinate $(p,E)$ use the finite state space
\[
 X_{p,E}=\{(a,d):0\leq a<p^{E-1},\quad1\leq d\leq p-1\}.
\]
Given an integer $x$, write $a=x\bmod p^{E-1}$, with the indicated
nonnegative representative. If $x\ne a$, write
$x-a=p^{E-1+t}u$, where $t\geq0$ and $p\nmid u$, and put
$d=u\bmod p$. If $x=a$, choose any $d$.
Thus every integer has a state in each coordinate. No assertion that this
state depends on a single finite modulus is needed.

A resource is an exponent vector $\mathbf e=(e_1,\ldots,e_s)$ with
$0\leq e_i\leq E_i$. Its smallest base modulus is
$b(\mathbf e)=\prod_i p_i^{e_i}$. A row assigns one label in every
coordinate. The coordinate condition is:
\begin{itemize}
\item $e=0$: no restriction;
\item $1\leq e<E$: $a\equiv r\pmod {p^e}$ for a chosen $0\leq r<p^e$;
\item $e=E$: one chosen pair $(a,d)\in X_{p,E}$.
\end{itemize}
The row is the product of these coordinate conditions. Select at most one
row for each exponent vector, and only use vectors with $b(\mathbf e)\geq M$.
The finite search asks whether the selected rows cover
$X=\prod_i X_{p_i,E_i}$.

The last condition represents the infinite family
\[
 x\equiv a+d p^{k-1}\pmod {p^k},\qquad k\geq E.
\]
It is essential to close this tail by a finite construction. An infinite
cover, or coverage of only sampled states, is not a certificate.

\section{Finite closure and distinctness}
\begin{lemma}
Fix $p,E,a,d$ as above, and a prime $q\ne p$. Put $N=E+q-1$.
There is a finite collection of congruence classes, with distinct moduli all
at least $p^E$, covering every integer whose state is $(a,d)$, and also $a$.
Every modulus has $p$-adic exponent at least $E$ and no other prime factor
except possibly $q$.
\end{lemma}
\begin{proof}
Use the $q$ classes
\[
 x\equiv a+d p^{k-1}\pmod {p^k},\qquad E\leq k\leq N,
\]
and the $q$ classes, defined by the Chinese remainder theorem,
\[
 x\equiv a\pmod {p^{N-j}},\qquad x\equiv j\pmod q,
 \qquad 0\leq j<q.
\]
The first group covers all specified nonzero digits at levels at most $N$.
Any remaining integer under consideration is congruent to $a$ modulo $p^N$.
Choose $j=x\bmod q$; the corresponding class in the second group covers it.
The moduli in the first group are $p^E,\ldots,p^N$; those in the second
group are $qp^N,\ldots,qp^E$. These are distinct and have the stated bounds.
All classes remain inside the prefix $a\bmod p^{E-1}$, although they may
cover additional digit states. Such additional coverage is harmless.
\end{proof}

\begin{proposition}\label{prop:construction}
If the finite rows cover $X$, there is a finite distinct covering system
of $\Z$ with all moduli at least $M$.
\end{proposition}
\begin{proof}
Assign a different fresh prime $q$ to every occurrence of a tail coordinate
in every row, taking all such primes outside $\{p_1,\ldots,p_s\}$.
For fixed coordinates use their single congruence, and for tail coordinates
use the preceding lemma. Within each row take all Cartesian products of the
coordinate classes and combine them by CRT. All coordinate moduli in a row
are coprime, so this operation is valid.

Every integer has a tuple in $X$; some row contains that tuple. Every
coordinate of that row covers the integer, so one of the CRT classes does
as well. Each row generates finitely many classes, and there are finitely
many rows.

For distinctness, a modulus containing an auxiliary prime identifies its
row uniquely. Within that row, unique factorization identifies each coordinate
modulus, so the Cartesian products have distinct moduli. If a modulus has no
auxiliary prime, its exponent at $p_i$ is exactly $e_i$ when $e_i<E_i$, and
at least $E_i$ when $e_i=E_i$. These disjoint exponent ranges recover the
resource vector uniquely. Since no vector is used twice, distinct rows cannot
produce equal untagged moduli. A tagged modulus cannot equal an untagged one.
Finally all output exponents at base primes are at least the resource
exponents, so every modulus is at least $b(\mathbf e)\geq M$.
\end{proof}

In fact the least output modulus is $\min b(\mathbf e)$: choosing the
first class of each tail realizes that base modulus. The completion need
not be enumerated. With the specified auxiliary primes the number of classes
generated by a row is the product of $2q$ over its tail coordinates.

\paragraph{A failed kernel gives an uncovered progression.}
If a tuple $((a_i,d_i))_i$ misses every row, take its CRT lift
\[
 x\equiv a_i+d_i p_i^{E_i-1}\pmod {p_i^{E_i}},\qquad i=1,\ldots,s.
\]
Every integer in this progression has the specified state in every coordinate.
For a tail row, the classes at levels $k>E_i$ and all its closing classes
require $x\equiv a\pmod {p_i^{E_i}}$; they cannot hit the lift with its
nonzero final digit and matching prefix. A different prefix also cannot hit
it. The only possible matching tail class is the one at level $E_i$, which
requires precisely the row's state. Fixed coordinate conditions behave as in
the kernel. Since every row misses at least one coordinate, the whole CRT
progression misses the completed construction. Thus the extra coverage from
finite closure cannot repair a failed kernel. The checker records this
explicit progression as well as the finite-state counterexample.

\paragraph{Expressiveness and limitations.}
The present tails each follow a single prefix and a single nonzero digit;
they do not encode arbitrary interacting recursive constructions as compactly
as Nielsen's notation. Nevertheless, any given finite distinct covering
system is representable using only fixed coordinates if the base primes
include all its prime factors and the depths exceed all its exponents.
A failed search at fixed primes, depths, or resource cap is therefore only
a restricted failure. It is not a proof that $M=43$ is impossible.

\section{Exact capacity checks}
Under uniform counting measure on $X$, a row has mass
\[
 w(\mathbf e)=\prod_i L(p_i,E_i,e_i)^{-1},\qquad
 L(p,E,e)=\begin{cases}p^e&e<E,\\p^{E-1}(p-1)&e=E.\end{cases}
\]
Consequently $\sum_{\mathbf e}w(\mathbf e)\geq1$ is necessary for a
kernel cover. It is not sufficient, because rows may overlap.
Since
\[
 \sum_{e=0}^E\frac1{L(p,E,e)}=\frac p{p-1},
\]
when every $p_i^{E_i}\geq M$, the unrestricted resource capacity is
\[
 \prod_i\frac{p_i}{p_i-1}
 -\sum_{\substack{1\leq n<M\\\text{all prime factors of }n\text{ among }p_i}}
 \frac1n.
\]
For $M=43$ and base primes through 11 this is exactly
$20641/21600<1$. Independently, the same sum of reciprocals excludes any
distinct cover whose moduli are all at least 43 and have no prime factors
greater than 11. This is an elementary density obstruction, not a new bound
for unrestricted covering systems.

For primes through 13 the resource capacity becomes
$682691/561600>1$. With base moduli capped at 10,000 it is
$24371213/20592000>1$. Thus the density test alone no longer decides the
finite search; overlaps and placement are essential.

\section{Implementation and initial experiments}
\texttt{code/construction\_kernel.py} contains the standard-library checker.
It validates labels, exponent ranges, the minimum modulus, and resource
uniqueness. It then recursively checks the union of the finite boxes,
memoizing identical active-row sets. A failed check returns an explicit
uncovered state; reaching a verification limit returns \texttt{UNKNOWN}.
Only exhaustive verification can return \texttt{VERIFIED}.

\texttt{code/search\_construction\_kernel.py} uses OR-Tools CP-SAT to
assign the labels. It initially imposes coverage of sampled states, then
adds exact counterexamples returned by the independent checker. Feasibility
of the sampled constraints is never presented as a cover. Solver time limits
are recorded as incomplete results. Solver infeasibility, if encountered,
would concern only the recorded finite family and would not have an
independently checked unsatisfiability proof in this implementation.

The positive controls are:
\begin{center}
\begin{tabular}{rrrr}
\toprule
Minimum modulus & Base primes & Depths & Kernel states\\
\midrule
3 & $2,3$ & $2,2$ & 12\\
6 & $2,3,5$ & $4,2,2$ & 960\\
\bottomrule
\end{tabular}
\end{center}
Both pass the independent exact checker. The naive fresh-prime expansion
specified above produces respectively 1,366 and 234,295 congruence classes;
these are deliberately inefficient validation examples.

Tests include 240 randomly generated small kernels compared against direct
enumeration, 51 finite-tail checks over complete periods, coordinate projection
checks, and rejection of duplicated resources. These support the code; the
mathematical justification of large certificates is Proposition~\ref{prop:construction}.

The record-target experiments use primes through 13 and through 19, depths
$(6,4,3,2,2,2)$ and $(6,4,3,2,2,2,2,2)$, respectively, and minimum modulus 43.
The JSON files retain budgets, resource caps, sampled-constraint counts,
solver statuses, candidate labels when available, and uncovered witnesses.
Both exhausted their time budgets (180 and 240 seconds) without a verified
cover. The last candidates respectively miss the complete progressions
\[
 x\equiv87304082151\pmod{649296648000},
\]
and
\[
 x\equiv51779995805607655\pmod{67740469989192000}.
\]
These witnesses reject those candidates, not all possible row assignments.
A supplementary minimum-10 run also exhausted its 90-second budget without
a certificate. Search runtimes and these limited failures should not be
interpreted as evidence of optimality of the record 42.

\paragraph{Reproduction.}
From the repository root, the exact checks require ordinary Python 3:
\begin{verbatim}
python code/test_construction_kernel.py
python code/construction_kernel.py data/construction-kernel-control-6.json
\end{verbatim}
For searches install \texttt{ortools}, either normally or into
\texttt{tmp/construction-deps}; the latter was used in this session.
The full search commands and final statuses are recorded in the companion
experiment JSON files and the repository README.

\section{Owens's thesis: locating and reserving the 42 branch}
The downloaded thesis is stored as \texttt{literature/owens-2014.pdf}.
Printed pages 8--10 (PDF pages 14--16) contain the prime-7 construction.
Its third input is
\[
 7^\uparrow(\ldots,3(2,4,3^\uparrow(1,2)),\ldots).
\]
The first occurrence of the $2$ gives the unique class with modulus 42:
\[
 H=\{x:x\equiv0\pmod2,\ x\equiv1\pmod3,\ x\equiv3\pmod7\}
   =10\pmod{42}.
\]
The same symbolic input also produces moduli $6\cdot7^d$ for $d\geq2$;
those exceed 42 and need not be deleted. Replacing just the first class
requires covering the part of $H$ not already covered by the retained
construction. Covering all of $H$ is a sufficient, stronger objective.

\paragraph{An explicit neighboring exchange.}
The $4$ in the next input of the displayed $3$ is the class $80\pmod{84}$:
it satisfies $x\equiv0\pmod4$, $x\equiv2\pmod3$, and $x\equiv3\pmod7$.
Moving that class to $52\pmod{84}$ would cover half of $H$, but opens
$80\pmod{84}$. The other half is $10\pmod{84}$ and cannot be assigned
another modulus 84. Thus even the first binary split creates a coupled
repair problem. The later prime-43 and prime-47 constructions work on
different inputs modulo 3, so changing those inputs has downstream effects.

\paragraph{Reservation audit.}
In the symbolic construction, Sections 3.1--3.3 use every modulus
$2^a3^b5^c\geq42$. For $c=0$, these are the pure powers $2^a$ with $a\geq6$,
the pure powers $3^b$ with $b\geq4$, and the mixed powers above the cutoff.
For $c=1$, the used ranges are $(b=0,a\geq4)$, $(b=1,a\geq2)$,
and $b\geq2$. For $c=2$ they are $(b=0,a\geq1)$ and $b\geq1$.
For $c\geq3$ all $a,b\geq0$ occur, including the extra tails on printed
page 7. These are reservations of exponent families, not a verification of
all the residues in the full thesis.

For a modulus $2^a3^b5^c7^d$ with $d\geq1$, the unambiguous terms of
Section 3.4 reserve at least the following inner exponents:
\begin{center}
\begin{tabular}{ll}
\toprule
$c$ & Reserved $(a,b)$\\
\midrule
0 & $a\geq3$; or $b=1,\ a\in\{1,2\}$; or $b\geq2,\ a\leq1$; or $(2,2)$\\
1 & $b=0,\ a\geq1$; or $b=1,\ a\leq4$; or $b\geq2$\\
2 & all $a,b\geq0$\\
$\geq3$ & $b=0,\ a\leq2$; or $b\geq1,\ a\leq1$\\
\bottomrule
\end{tabular}
\end{center}
The extra $(2,2)$ is the $9\cdot4$ allocation on printed page 10, which
supports the prime-31 step. We deliberately omit the final $125$-tail
additions and the extra $125^\uparrow8^\uparrow$ from this reservation
table: the resulting spare pool is an \emph{overestimate}. In particular,
the calculation does not depend on resolving the compressed placement of
the final $125$ factors between the prose on page 8 and summary on page 9.

\begin{proposition}[A restricted density obstruction]
If these infinite symbolic families stay reserved, spare moduli with prime
factors only in $\{2,3,5,7\}$ and minimum modulus 43 cannot, by themselves,
cover the whole progression $H$.
\end{proposition}
\begin{proof}
A single class of modulus $n$ meets $H$ with relative density either zero
or $\gcd(n,42)/n$. The complement of the reservation table, after imposing
$n\geq43$, has the following total conditional capacities:
\begin{center}
\begin{tabular}{lr}
\toprule
Potential spare family & Sum of $\gcd(n,42)/n$\\
\midrule
$7^d,\ 2\cdot7^d,\ 3\cdot7^d$, $d\geq2$ & $1/2$\\
$4\cdot7^d$, $d\geq2$ & $1/12$\\
$5\cdot7^d$, $d\geq2$ & $1/30$\\
$4\cdot3^b7^d$, $b\geq3,d\geq1$ & $7/72$\\
$2^a3\cdot5\cdot7^d$, $a\geq5,d\geq1$ & $7/240$\\
$2^a5^c7^d$, $a\geq3,c\geq3,d\geq1$ & $7/1200$\\
$2^a3^b5^c7^d$, $a\geq2,b\geq1,c\geq3,d\geq1$ & $7/400$\\
\midrule
Total & $2759/3600<1$\\
\bottomrule
\end{tabular}
\end{center}
Each entry is a convergent geometric sum. The union bound for any finite
selection is strictly less than one, which proves the claim.
\end{proof}

Even releasing all the moduli $36\cdot7^d$ allocated to the prime-31 step
raises this bound only to $1153/1200<1$. This does \emph{not} rule out
covering only the true residual subset of $H$, using later-stage overlaps,
rearranging occupied families, or releasing high powers after a particular
finite truncation. In particular it is not a proof that a local repair of
Owens's cover is impossible.

\section{Fresh-prime partial replacements and their limitation}
There is a constructive route that certifies part of the branch without
needing a complete transcription of Owens's later stages. All his symbolic
moduli have prime factors at most 89. A new largest prime $q\geq97$ can
therefore distinguish a new family of moduli from the old ones.

\begin{lemma}[Disjoint set-building rule]\label{lem:sets}
Let $F=a\pmod m$. Suppose $T$ symbolic sets each cover $F$ and have
pairwise disjoint modulus families. Let $p$ be a prime absent from these
families, and write $e=v_p(m)$. One can construct
\[
 T'=(e+1)T+\left\lfloor\frac{T}{p-1}\right\rfloor
\]
such sets using the additional prime $p$.
\end{lemma}
\begin{proof}
Keep the $T$ old sets. For each $j=1,\ldots,e$, take a copy of every old
set with the extra condition $x\equiv a\pmod{p^j}$; its moduli are
multiplied by $p^j$. These copies still cover $F$.
Divide the old list into disjoint groups of $p-1$ sets, discarding an
incomplete final group. Each group fills one $p^{e+1}$-arrow rooted at
$a\pmod{p^e}$. It covers $F$ symbolically, with the usual finite closure
at the exceptional tail. Distinctness follows from the disjoint exponent
ranges $0,1,\ldots,e,[e+1,\infty)$ at $p$, and from disjointness of
the old families within and between the groups.
\end{proof}

Start with the formal unit set, with modulus 1, and process each chosen
prime once by this rule. If the final count is at least $q-1$, use any
$q-1$ sets as the inputs of $q^\uparrow$. Every resulting modulus is at
least $q$. Finitely many nested arrow occurrences can be closed with
distinct auxiliary primes outside the base primes, using the finite-tail
construction above. This gives an actual finite partial covering of $F$.
It need not be expanded into an enormous list to specify the construction.

\paragraph{Six certified partial replacements.}
Take $m=840=2^3\cdot3\cdot5\cdot7$. The following master order was found
by permutation search:
\begin{align*}
(&2,5,11,3,7,61,67,17,59,19,37,79,73,41,71,83,\\
 &43,29,23,31,47,13,103,53,97,109,89,107,101).
\end{align*}
For each row below, remove all primes at least $q$ from this order,
start with $T=1$, and apply Lemma~\ref{lem:sets} with
$(e_2,e_3,e_5,e_7)=(3,1,1,1)$ and all other exponents zero.
\begin{center}
\begin{tabular}{rrrr}
\toprule
Progression covered & Fresh root $q$ & Available sets & Required sets\\
\midrule
$10\pmod{840}$ & 97 & 109 & 96\\
$52\pmod{840}$ & 101 & 110 & 100\\
$94\pmod{840}$ & 103 & 111 & 102\\
$136\pmod{840}$ & 107 & 112 & 106\\
$178\pmod{840}$ & 109 & 113 & 108\\
$220\pmod{840}$ & 113 & 114 & 112\\
\bottomrule
\end{tabular}
\end{center}
For example, the $q=97$ recurrence starts
$1\to5\to11\to12\to30\to65$ on primes $2,5,11,3,7$.
All six progressions lie in $H$ and are disjoint. Together they cover
$6/20=3/10$ of $H$. This is a fraction of the \emph{whole} branch,
not a measured fraction of the actual gap after deleting 42; some of
these points can already be covered elsewhere in Owens's construction.
The new partial covers alone certainly do not cover all of $H$: their
stated completions all miss
\[
 x\equiv189833261760964\pmod{1117116169368360}.
\]
This witness is $724\pmod{840}$ and has residues
$(4,6,4,4,4,6)$ at the six root primes. Those nonzero root digits select
either the fixed-8 or fixed-5 unit child, whose residue it fails.
It is not asserted to miss Owens's retained classes.

The six modulus families are disjoint: the largest nonauxiliary prime in
each identifies its root $q$. An earlier root never uses a later prime.
Owens's original families have no prime above 89. Auxiliary primes for all
finite closures must be selected jointly, outside every listed prime and
distinct across occurrences, so they do not invalidate this separation.
This establishes finite partial covers with no modulus 42; it does not
establish a covering system with minimum modulus 43.

\paragraph{Why this rule stalls.}
For a fixed progression and set of base primes, every order in this rule
satisfies the rational relaxation
\[
 T\leq B=\prod_{p}\left(e_p+1+\frac1{p-1}\right).
\]
For one class modulo 840 and \emph{all} primes below 127,
$B=121.360689\ldots<126$. Hence no order of these operations can fill
a $127^\uparrow$ on such a class. Nor can a larger fresh root help within
this rule: if $B(q)<q-1$, adding the base prime $q$ gives
$B(q)q/(q-1)<q$, whereas the next prime requires at least $q+1$ inputs.
This induction concerns the stated counting rule only; it is not an upper
bound on all possible partial coverings of a progression.

We also searched 72 uniform split shapes
\[
 m=2^r3^s5^t7^u,\quad
 1\leq r\leq6,\ 1\leq s\leq2,\ 0\leq t\leq2,\ 1\leq u\leq2.
\]
The number of equal pieces of $H$ is
$K=2^{r-1}3^{s-1}5^t7^{u-1}$. With base primes frozen at 89 and one
different fresh root for each piece, the rational relaxation is below
the inputs required by the $K$th fresh prime in all 72 cases.
Those 72 restricted uniform constructions are excluded independent of
the permutation search. Nonuniform splits, partial inputs, and borrowing
families from the original cover remain outside this calculation.

\paragraph{Reproducibility and remaining work.}
\texttt{code/owens\_branch\_search.py} writes the exact rational sums,
all 72 configurations, and the six partial construction schemas to
\texttt{data/owens-42-branch-search.json}. Run
\begin{verbatim}
python code/owens_branch_search.py --iterations 2500
python code/test_owens_branch_search.py
\end{verbatim}
The checker validates the construction recurrence, progression containment,
prime separation, and available input counts. It does not claim to have
enumerated the finite completion or verified Owens's full cover. The
geometric-sum audit is also cross-checked against 7,280 exponent tuples.
Twelve small literal expansions separately pass exhaustive digit-state
coverage checks and pairwise exponent-family intersection checks; the
six actual partial schemas pass the count and progression checker.

The next coupled instance should include the $42$ and $84$ branches,
the $36\cdot7^d$ allocation that supports prime 31, and the early
prime-3 placement. The extra $81^\uparrow$ is another candidate for
relocation, as Owens himself suggests on printed pages 18--19. Before a
successful switch can be asserted, its opened holes and the later stages
that rely on their present coverage must all be encoded and checked.
That full residual-cover computation remains unfinished here. No
minimum-modulus-43 covering system has been obtained.

\section{The new construction target: minimum modulus at least 1000}
The user's new target is a \emph{construction lower bound}: a finite
cover with distinct moduli, all at least 1000. It is not a request to
lower a universal upper bound to 1000. No such construction was obtained.
The contemporary primary literature inspected for this continuation still
reports Owens's minimum modulus 42~\cite{agrawal}; this limited source
check is not an exhaustive literature survey.

\paragraph{The scale of the change to Owens's construction.}
The reservation audit above identifies 64 moduli involving only $2,3,5$
and 42 further moduli involving $7$, all in $[42,999]$, already allocated
in the early stages. Thus at least 106 early classes would need to be
removed or replaced to reach 1000 while starting from that schedule.
The full list is saved in \texttt{data/construction-1000-profile.json}.
This count does not include the many later-prime classes below 1000.

\paragraph{Exact necessary conditions.}
For a set of primes $\mathcal P$, define
\[
 C(M,\mathcal P)=
 \prod_{p\in\mathcal P}\frac{p}{p-1}
 -\sum_{\substack{n<M\\p\mid n\Rightarrow p\in\mathcal P}}\frac1n.
\]
Unique factorization and the geometric series show that this is exactly
the reciprocal sum over all $\mathcal P$-smooth integers at least $M$.
Every finite distinct cover using only these primes must have
$C(M,\mathcal P)\geq1$. With all primes at most $y$ allowed, the exact
rational computations give the following decimal displays:
\begin{center}
\begin{tabular}{rr}
\toprule
$y$ & $C(1000,\{p:p\leq y\})$\\
\midrule
7 & 0.0844650041\\
47 & 0.9116749285\\
53 & 0.9843755046\\
59 & 1.0537687824\\
\bottomrule
\end{tabular}
\end{center}
Consequently a minimum-1000 cover must use a prime divisor at least 59.
Passing the test at 59 is not an existence result. The very small mass
available from the first four primes also explains why retaining the
early small-prime layout supplies little coverage after the new cutoff.

Separately, if there are $k$ distinct moduli, their reciprocal sum is at
most $\sum_{j=0}^{k-1}(1000+j)^{-1}$. Exact rational sums show
\[
 \sum_{n=1000}^{2716}\frac1n<1
 <\sum_{n=1000}^{2717}\frac1n.
\]
Thus at least 1718 classes are necessary, and the largest modulus must
be at least 2717. These are elementary necessary conditions, not sharp
predictions of the size of a successful construction.

\paragraph{Kernel screening.}
We use the existing finite-tail model with
$E_p=\min\{e\geq1:p^e\geq1000\}$. Recursive generation now prunes a
partial exponent vector as soon as its base product exceeds the cap;
it avoids enumerating the full product of exponent choices. Several
representative finite families are:
\begin{center}
\begin{tabular}{rrrr}
\toprule
Largest base prime & Base-product cap & Resources & Kernel capacity\\
\midrule
59 & 10000 & 2189 & 0.659943\\
73 & 30000 & 6142 & 0.986637\\
97 & 30000 & 7210 & 1.127651\\
149 & 10000 & 3801 & 1.083089\\
311 & 10000 & 4995 & 1.388128\\
311 & 30000 & 13212 & 1.851013\\
\bottomrule
\end{tabular}
\end{center}
The first two are excluded by exact capacity. These caps refer to
\emph{resource bases}, not the largest moduli in a finite completion:
tail exponents and auxiliary primes can make completed moduli much larger.

\paragraph{Exact verification at larger prime supports.}
The verifier now partitions each coordinate with a sparse tree of its
base-$p$ digits. Fixed residue conditions end at prefix nodes; tail
conditions end at the final digit. Every unmentioned child has the same
active rows, so one representative suffices. This preserves the exhaustive
coverage test while avoiding a scan of every state in a coordinate.
Against direct enumeration, 160 coordinate-partition comparisons and
240 complete small-kernel comparisons pass. The checks also include
30 capped-resource comparisons, 51 finite-tail tests, and an explicit
coordinate with $2^{49}$ states handled symbolically.

\paragraph{Bounded searches performed.}
One CP-SAT run used the prime-149 family, 3801 resources, and a
180-second budget. It found a candidate satisfying 65 sampled constraints.
The exact checker rejected that candidate. After adding constraints to
reach 130 sampled points, the remaining solves returned unknown before
the time budget expired. This is not an infeasibility result.

A second search chooses each resource's residue tuple by the gain in
covered training points, then revisits resources in randomized order.
Training codes are cached for small instances and streamed for larger
ones. Fresh diagnostic points are not used to choose the best candidate.
The results below all use base-product cap 10000; each validation sample
has 32768 points.
\begin{center}
\begin{tabular}{rrrrr}
\toprule
Largest prime & Training points & Sweeps & Training misses & Validation misses\\
\midrule
149 & 32768 & 1 & 0 & 11572\\
149 & 262144 & 6 & 38411 & 10905\\
311 & 262144 & 6 & 15693 & 7944\\
\bottomrule
\end{tabular}
\end{center}
The entries are for the best training candidate in each run. The second
run increases the weights of uncovered points after each sweep; the
third keeps uniform weights. Seeds are 1000, 1001, and 1002 respectively.
All three candidates fail exact verification. In particular, even perfect
coverage of the first training sample was a misleading indicator of
global coverage. The sample counts are not certified density estimates
and do not establish proximity to a covering system.

Each rejected candidate has an explicit uncovered arithmetic progression
in its JSON record. The separate standard-library program
\texttt{verify\_construction\_searches.py} checks these witnesses without
using the coordinate partition algorithm: for every row it exhibits a
prime-power condition incompatible with the progression, including the
whole tail and its finite closing classes. A small positive control at
minimum modulus 3 passes exhaustive kernel verification. An additional
minimum-10 diagnostic candidate is rejected, despite fitting its training
sample; it is not reported as a positive control.

\paragraph{Reproduction.}
The capacity profile and large searches are reproducible with:
\begin{verbatim}
python code/search_large_minimum.py --profile 
  --output data/construction-1000-profile.json
python code/search_large_minimum.py --largest-prime 149
  --samples 262144 --sweeps 6 --seconds 180 --seed 1001
  --output data/construction-greedy-1000-149-large.json
python code/search_large_minimum.py --largest-prime 311
  --samples 262144 --sweeps 6 --seconds 180 --seed 1002
  --uniform-weights --output data/construction-greedy-1000-311.json
\end{verbatim}
Join each displayed command's wrapped lines when running it. NumPy is
needed for the sample-guided search; capacity screening and witness
verification use only the standard library. The earlier CP-SAT script
still requires OR-Tools. The rational values, candidate assignments, seeds,
budgets, diagnostics, and witnesses are retained in the data files.

The work so far gives a scalable search and reliable rejection procedure,
not the requested construction. The present random-sample objective
substantially overfits even at 262144 points. A useful continuation must
incorporate structural residual constraints or a different recursive
construction, rather than treating sample coverage as the success test.
Neither minimum modulus 1000 nor minimum modulus 43 has been established.

\section{Target 100: exact overlap control and fresh-prime blocks}
\label{sec:structural100}

This continuation changes the construction rule rather than increasing
the training sample. No sampled points are used. Its positive control has
minimum modulus 10; its target-100 constructions remain incomplete.

\subsection{An exact residual partition}
Write each coordinate state as its $E_p-1$ ordinary base-$p$ digits followed
by its first nonzero tail digit, whose alphabet has size $p-1$. A row of
exponent $e<E_p$ fixes the first $e$ ordinary digits. A tail row fixes all
$E_p$ digits. Thus rows are cylinders in a finite product of digit alphabets.
Represent the residual set $R$ by a reduced ordered decision diagram.
For a resource fixing the digit positions $S$, compute
\[
 g_S(u)=\#\{x\in R:x|_S=u\}
\]
by summing over all other positions. Place the resource at a maximizing
$u$ and replace $R$ by $R\setminus\{x:x|_S=u\}$. This calculation counts
every overlap exactly. A second placement rule chooses digits successively
using exact marginal counts; it is a heuristic for the joint maximization,
but its final gain and the resulting residual are still exact.

The search uses integer arithmetic throughout. A diagram limit stops a
search with its last valid partial construction, and does not prove
infeasibility. A separate verifier, based on coordinate cylinders rather
than marginal decision diagrams, checks any proposed cover.

The minimum-10 control found 215 exponent resources on primes through 23.
The independent verifier certified all $91\,968\,307\,200$ kernel states,
using 272272 verification nodes. The finite-completion rule gives an
explicit schema with $171\,744\,802$ distinct classes, smallest modulus 10,
and largest auxiliary prime 1483. This is a correctness control, not an
efficient construction or a record improvement.

For the target 100, use primes through 31 and depths
\[
 (7,5,3,3,2,2,2,2,2,2,2).
\]
With resource base at most 1000, there are 348 available resources.
The digitwise placement run stops at its 200000-node diagram limit after
placing 257 rows. Its exact uncovered kernel fraction is
\[
 r_0=\frac{548023176592474367}{1398918654701568000}
     =0.391747707953315\ldots.
\]
Other exact runs with different resource orders and coordinate orders also
stop short of coverage. All stored target-100 base candidates have explicit
uncovered progressions checked independently by modular arithmetic.

\subsection{A structural repair with a separate prime for each block}
Let $q$ be a fresh prime, greater than every base prime. Choose $q-1$
distinct positive integers $n_1,\ldots,n_{q-1}$, supported on base primes,
with $qn_d\ge100$. Set $m=\operatorname{lcm}(n_1,\ldots,n_{q-1})$ and
choose a residue $a\pmod m$. Use the families
\begin{equation}\label{eq:divisor-block}
 x\equiv a\pmod {n_d},\qquad
 x\equiv d q^{k-1}\pmod {q^k},
 \quad 1\le d\le q-1,\quad k\ge1.
\end{equation}
The resulting moduli $n_dq^k$ are distinct: $q\nmid n_d$, and the $n_d$
are distinct. Blocks assigned different fresh primes cannot collide,
because their $n_d$ use only the original base primes. They also cannot
collide with the base construction. Finite completion uses a different
auxiliary prime for each tail occurrence, all beyond the structural primes.
Every modulus remains at least 100.

\begin{proposition}[Exact projection identity]
Suppose each $n_d$ uses only fixed prefix exponents $e_p<E_p$. For a
fixed base kernel state $x$, block~\eqref{eq:divisor-block} covers every
state of its new $q$ coordinate if and only if $x\equiv a\pmod m$.
For several blocks with different fresh roots and no dependencies on one
another, a base state has an uncovered extension to the new coordinates
if and only if it is uncovered by the base and lies outside every one of
these progressions.
\end{proposition}
\begin{proof}
At root state $d$, the block covers $x$ exactly when $x\equiv a\pmod{n_d}$.
All these congruences hold exactly when $x\equiv a\pmod m$. If no block
covers all its root states, choose a missed state separately at each root.
Their Cartesian product is an uncovered extension. Conversely, a block
covering every state of its root covers every extension of $x$.
\end{proof}

This identity controls the overlap between blocks exactly, including
overlap with the base. The construction searches for prefix cylinders
having maximum intersection with the current projected residual, and
erases only that intersection. It emits the actual $q-1$ kernel rows for
each block, rather than treating a virtual cylinder as one modulus.

The implemented experiment considers prefix moduli supported on $2,3,5,7$.
For each root, it compares the first four distinct admissible moduli,
optimizing the residue of each by the exact marginal calculation.
Nineteen blocks with roots $37,41,\ldots,113$ add 1414 rows to the base's
257 rows. For example, the first block has $q=37$, $a=1323$, $m=1680$;
its divisors are the first 36 divisors of 1680 satisfying $37n\ge100$.
The last block has $q=113$, $a=31500$, $m=60480$.
The final fraction of base states admitting an uncovered extension is
\[
 r_1=\frac{11122926985118327}{28549360300032000}
     =0.389603370030882\ldots.
\]
This is a projected kernel fraction. It is not the uncovered natural
density of the finite completion, nor the uncovered fraction of the
enlarged kernel. All block residues, divisor assignments, and exact
incremental gains are saved with the emitted rows.

An uncovered projected state can be lifted efficiently: at each root
choose a divisor $n_d$ whose congruence it misses, assign that root the
nonzero residue $d\pmod q$, and use CRT. The independent checker then
tests the resulting progression against every actual row. Later root-tail
classes and closing classes require $x\equiv0\pmod q$, whereas the witness
has $x\equiv d\not\equiv0\pmod q$. The base rows are checked against their
full finite completions in the same way as before. Thus this procedure
can certify failure without exhausting the general union verifier.
For the saved 1671-row candidate the independent check succeeds: every
integer $x\equiv R\pmod L$ is missed, where
\begin{align*}
 R&=29370484486033522038888368677610496702136998137198482290242400,\\
 L&=191713376283959466494384635189844768308141224115446848519088000.
\end{align*}
The row-by-row modular check includes the finite closing classes, so this
is a genuine failure certificate for the finite-completion schema.

\subsection{A capacity obstruction for the whole independent-block repair}
The weak gain above is structural. Freeze the 257-row base and allow
\emph{all} its base primes in the prefix moduli, with $e_p<E_p$.
Set $B=\prod_p p^{E_p-1}$ and define
\[
 m_*(q)=\min\left\{m:m\mid B,\quad
       \#\{n\mid m:qn\ge100\}\ge q-1\right\},
\]
omitting roots with no such $m$. There are
$\prod_p E_p=40320$ possible prefix moduli, and at most 40320 divisors
in any one. Hence no root above 40321 is admissible.
Every block at root $q$ removes a projected cylinder of mass at most
$1/m_*(q)$. The union bound therefore gives the following necessary
condition for a complete repair using at most one such block per root:
\[
 r_0\ \le\ \sum_{\substack{q>31\\q\text{ prime}}}\frac1{m_*(q)}.
\]
An exact enumeration of all 40320 prefix shapes and all 4220 admissible
roots gives
\[
 \sum_q\frac1{m_*(q)}
 =\frac{9001713966181}{3032474610765600}
 =0.002968438361932\ldots <r_0.
\]
Consequently, even using every admissible root and optimizing all
residues, this restricted repair leaves projected mass at least
\[
 \frac{488939644940579871613}{1257627870576709632000}
 =0.388779269591383\ldots.
\]
The CRT argument above turns a remaining state into an uncovered
progression for the associated finite completions. This is a rigorous
exclusion for the frozen-base independent-divisor-block method. It does
not exclude adding unused base resources, changing the base, allowing
several interacting blocks at one root, using earlier root primes inside
later blocks, or using more general complete-set inputs.

\subsection{Verification and reproducibility}
The new tests compare 900 exact marginal and residual calculations with
literal enumeration, and test row encoding in all three coordinate
orders. An additional 2880 exhaustive small-state checks verify the
fresh-prime projection identity; 14 small capacity audits compare the
divisor-count screen with direct enumeration. The earlier kernel tests,
including 51 finite-tail checks, also pass.

The main scripts are \path{code/structural_cover_search.py} and
\path{code/structural_prime_patches.py}. The base, repair, and capacity
artifacts are respectively
\begin{itemize}
\item \path{data/construction-structural-100-31-digits.json};
\item \path{data/construction-structural-100-prime-patches.json};
\item \path{data/construction-structural-100-patch-capacity.json}.
\end{itemize}
The independently verified positive control is
\path{data/construction-structural-control-10-verified.json}.
Run the standard-library tests with
\begin{verbatim}
python code/test_structural_cover_search.py
python code/test_structural_prime_patches.py
python code/test_construction_kernel.py
\end{verbatim}
The accompanying README contains the full search commands. The next
structural change must overcome the independent-block projection
obstruction, for example by arranging dependencies between root primes
and jointly reallocating base resources. Simply adding more roots within
this fixed-divisor repair cannot work. No minimum-100 cover, or
improvement over 42, is claimed.

\section{Target 60: capacity screening and a disjoint-cylinder construction}
\label{sec:structural60}

Lowering the target from 100 to 60 relaxes the density constraints, but
the searches below do not produce a minimum-60 covering system.
For primes $P$, the reciprocal sum of every $P$-smooth modulus at least 60 is
\[
 C(60,P)=\prod_{p\in P}\frac{p}{p-1}
       -\sum_{\substack{n<60\\n\text{ is }P\text{-smooth}}}\frac1n.
\]
For $P$ consisting of all primes through 11 this equals
$3084667/3880800=0.7948533807\ldots<1$. Including 13 raises it to
$104496767/100900800=1.0356386372\ldots$. Hence any distinct cover with
minimum at least 60 must use a prime divisor at least 13; this screen
does not establish that primes through 13 suffice. The corresponding
first nonexcluded prime range for target 100 ends at 19.
Also, the reciprocal sum of the 102 smallest permitted distinct moduli,
$60,\ldots,161$, is less than one, while adding 162 exceeds one.
Thus at least 103 classes are necessary, as is a maximum modulus at least 162.

The exact searches use the same kernel formulation and independent
verifier as before. The depths through prime 31 are now
$(6,4,3,3,2,2,2,2,2,2,2)$. With resource base at most 1000, 369
resources are available, with total kernel capacity
$10720444249609/7487591631520=1.4317613429\ldots$.
Through prime 43 with base cap 3000 there are 973 resources and capacity
$1.9635285411\ldots$. Capacity exceeding one does not control overlap.

\begin{center}
\begin{tabular}{lrrl}
\toprule
Placement / coordinate order & Rows & Uncovered fraction & Stop reason\\
\midrule
Digitwise / increasing primes & 260 & 0.27185609 & Node limit\\
Maximum / decreasing primes & 257 & 0.30525110 & Node limit\\
Digitwise / decreasing primes & 245 & 0.39283632 & Node limit\\
Mass-ordered / decreasing primes & 59 & 0.47534648 & Node limit\\
Disjoint cylinders / decreasing primes & 180 & 0.74502687 & Resources exhausted\\
\bottomrule
\end{tabular}
\end{center}
The first run uses primes through 31, cap 1000, and 200000 diagram nodes.
The next three use primes through 43, cap 3000, and 100000 nodes.
The last uses primes through 43, cap 10000, and 100000 nodes.
All are deterministic and use exact integer counts; no samples are drawn.
The listed fractions refer to each finite kernel, not to the natural
density of integers missed by finite completions. These different search
budgets and kernel spaces also preclude interpreting the table as a
controlled ranking of the placement rules.

\paragraph{The additional structural policy.}
For selected digit positions $S$, instead of maximizing a marginal count,
form the Boolean function
\[
 h_S(u)=\bigwedge_{v\text{ on the other positions}}1_R(u,v).
\]
If some $u$ has $h_S(u)=1$, its whole cylinder lies in the residual;
choose such a cylinder and erase it. Otherwise skip this resource.
Universal quantification is implemented by Boolean intersections in a
reduced decision diagram. Every placed cylinder is disjoint from all
previously placed cylinders. This restriction is on the kernel only:
the finite closing classes may overlap.
The run inspected all 2098 permitted resources, placed 180 cylinders,
and stopped with a nonempty residual. Independent pairwise checks verify
all $\binom{180}{2}=16110$ disjointness claims, and the sum of cylinder
masses agrees exactly with one minus the residual fraction. The universal
quantification calculation also passes 900 comparisons with literal
enumeration. Failure of this greedy run does not exclude other disjoint
assignments or other resource pools.

\paragraph{A concrete failure certificate.}
The strongest partial result in this batch leaves the exact fraction
\[
 \frac{1377915076978579}{5068545850368000}
 =0.271856093967964\ldots
\]
uncovered. Every integer in
\[
 x\equiv65199869485917100379505775
 \pmod{202731531413976570181176000}
\]
misses its 260 rows and their finite completions. The modular checker
independently verifies this statement. It likewise verifies explicit
uncovered progressions for each of the other four candidates.

For this frozen 260-row base, the earlier independent-divisor-block
capacity obstruction still applies. Allowing all prefix moduli on its
original prime coordinates gives 27648 prefix shapes and 3005 admissible
fresh roots, with total removable projected mass at most
\[
 \frac{163159827287}{53201308960800}=0.003066838588637\ldots.
\]
This excludes completing that base by at most one independent fixed-divisor
block per root. It does not exclude changing the base or introducing
dependencies between blocks. No additional patch experiment is needed
to reject that restricted route.

\paragraph{Saved evidence.}
The profile, witness checks, and disjointness checks are
\path{data/construction-60-profile.json},
\path{data/construction-60-witness-checks.json}, and
\path{data/construction-60-disjoint-check.json}.
The best partial assignment is
\path{data/construction-structural-60-31-digits.json}.
The restricted repair audit is
\path{data/construction-structural-60-patch-capacity.json}.
Reproduction commands are in the README. The present evidence establishes
neither existence nor impossibility of a minimum-60 distinct covering system.

\section{Squarefree moduli: target 30 and the literature barrier}
\label{sec:squarefree30}

Here every modulus must be squarefree, the moduli must be pairwise distinct,
and the system must be a finite cover of all integers. The literature checked
does not supply an example with minimum modulus greater than 2.
Dalton and Trifonov~\cite{daltontrifonov}, Problem 4, ask whether every
distinct squarefree covering must have least modulus 2. An example with
minimum at least 30 would answer that question negatively. Any example
with minimum greater than 3 would also yield one with minimum exactly 3
by adding an arbitrary class modulo 3.
Cummings, Filaseta, and Trifonov~\cite{cft} prove the upper bound 118 for
the minimum of a distinct squarefree covering. That upper bound does not
settle existence at 30. No impossibility at 30 is proved here.

\paragraph{The correct finite model.}
An exponent capped at $E_p=1$ in the preceding tail model represents
\emph{all} powers $p^k$, $k\ge1$, rather than a squarefree modulus.
Its finite completion still contains higher powers of $p$. Thus changing
the earlier depth parameter to 1 would not impose squarefreeness.
Instead use the ordinary CRT box $\prod_{p\in P}\Z/p\Z$, with $p$ states
at prime $p$, and at most one row for each nonempty subset $S\subseteq P$,
of modulus $\prod_{p\in S}p$. There are no tail rows or auxiliary-prime
completions. Every output is a literal finite congruence list.

\paragraph{An exact necessary prime-support condition.}
Writing $Q=\prod_{p\le y}p$, the total reciprocal capacity is
\[
 C_{\rm sf}(30,y)=\sum_{\substack{n\mid Q\\n\ge30}}\frac1n
 =\prod_{p\le y}\left(1+\frac1p\right)
  -\sum_{\substack{n\mid Q\\n<30}}\frac1n.
\]
Exact arithmetic gives
\[
 C_{\rm sf}(30,31)=\frac{31160796824}{33426748355}<1,
 \qquad
 C_{\rm sf}(30,37)=\frac{5822269528}{5596333435}>1.
\]
Consequently, any squarefree distinct covering with minimum at least 30
must have a modulus divisible by a prime at least 37. The second inequality
does not prove that primes through 37 suffice.

\paragraph{The bounded experiment and its verification.}
The separate script \path{code/search_squarefree_construction.py} uses
exact marginal placement in the ordinary CRT box. Through prime 37 there
are 4078 eligible squarefree resources and period $Q=7420738134810$.
The run places 121 rows before reaching its 100000-node diagram limit.
Its exact uncovered natural density is
\[
 \frac{52284300429}{145504669310}=0.359330739535289\ldots.
\]
Unlike earlier tail-kernel fractions, this is the actual density of missed
integers for the emitted finite congruence list. Direct modular arithmetic
verifies that every integer in
\[
 x\equiv5045898743977\pmod{7420738134810}
\]
misses all 121 classes. Each modulus divides the squarefree period $Q$,
and all are distinct and at least 30.

For the independent union checker only, ordinary residues are encoded
as prefix rows with $E_p=2$ and exponents $e_p\in\{0,1\}$. Every used
exponent is strictly below its tail threshold. The extra coordinate states
are unused duplicates and do not change coverage. Two hundred small
direct-period comparisons check this encoding and the verifier.
A positive control produces the following 14-class squarefree cover:
\[
\begin{array}{c|rrrrrrr}
n&2&3&5&6&7&10&14\\
a&0&0&0&1&0&1&1\\ \hline
n&15&21&30&35&42&70&105\\
a&2&2&23&24&11&19&104
\end{array}
\]
Every residue modulo 210 was checked directly; its least modulus is 2.
It is a correctness control, not a new record.

The failed target-30 candidate, positive control, and test results are
\path{data/construction-squarefree-30-37.json},
\path{data/construction-squarefree-control-2.json}, and
\path{data/construction-squarefree-tests.json}. The run's failure does not
exclude other placements on these primes, or larger prime supports.

\section{Squarefree minimum 3: normalized parity layers and joint repair}
\label{sec:parity3}

The new search coordinates the odd classes and the even classes across
both parity layers. Every modulus is a literal squarefree integer; modulus
2 is forbidden. No sampling or prime-power completion is used. No complete
minimum-3 covering has been found.

\subsection{A lossless normalization}
Fix distinct odd primes $p_1,\ldots,p_k$, with $p_1=3$, and period
$L=2\prod_i p_i$. For existence, we may add every unused divisor of $L$
other than 1 and 2 as a new modulus, choosing arbitrary residues. Adding
classes cannot destroy coverage. In particular, classes modulo each $p_i$
and $2p_i$ can be assumed present.

\begin{lemma}[Prime and twice-prime normalization]
If a squarefree distinct covering exists on this prime support with all
moduli at least 3, then one exists with the following properties. The class
modulo $p_i$ fixes its CRT coordinate to 0. The class modulo $2p_i$ fixes
that coordinate to 1 and its parity to some $b_i\in\{0,1\}$. We may take
$b_1=0$. Every other class has nonzero labels at its fixed odd coordinates.
\end{lemma}
\begin{proof}
If the classes modulo $p_i$ and $2p_i$ have the same label at $p_i$, the
latter class is redundant. Move its $p_i$ label to any different value;
this loses no coverage because the former class remains. A permutation
of the $p_i$ coordinate can then send these two labels to 0 and 1.
Apply this permutation to every class using $p_i$. Independent coordinate
permutations give a bijection of the CRT box and preserve all cylinder
supports, hence all moduli, distinctness, and coverage. Flip the parity
coordinate if necessary to make $b_1=0$. Finally, any composite class
having a zero odd label lies inside an already present prime class.
It can be reassigned nonzero odd labels without losing coverage.
\end{proof}

This normalization also cannot increase the number of uncovered points,
starting from any partial system. Thus lower bounds for the normalized
systems apply to all systems on the same prime support.
Once the prime classes are placed, the only unhandled points are the two
copies of $\prod_i\{1,\ldots,p_i-1\}$. In layer $\epsilon$, placing the
classes modulo $2p_i$ leaves a product box with coordinate sizes
\[
 D_i^\epsilon=p_i-1-\mathbf1_{b_i=\epsilon},\qquad
 R_\epsilon=\prod_iD_i^\epsilon.
\]
The solver chooses the parity labels and all composite classes jointly,
with constraints for every point in this finite box. It minimizes the
number left uncovered. Solver optimality is not assumed for the results below.

\subsection{A checkable obstruction through prime 11}
For each subset $S\subseteq\{1,\ldots,k\}$ with $|S|\ge2$, there are two
remaining resources: $d_S=\prod_{i\in S}p_i$ and $2d_S$. Define
\[
 V_\epsilon(S)=\prod_{i\notin S}D_i^\epsilon.
\]
The odd class covers at most $V_0(S)+V_1(S)$ residual points. The even
class occupies only one parity layer and covers at most
$\max\{V_0(S),V_1(S)\}$. Therefore the number of missed points is at least
\begin{equation}\label{eq:anchor-capacity}
 R_0+R_1-\sum_{|S|\ge2}\bigl(V_0(S)+V_1(S)+\max(V_0(S),V_1(S))\bigr).
\end{equation}
This is a union bound; overlaps can only make the uncovered set larger.
For $p_i=3,5,7,11$, all parity patterns up to flipping both layers give:
\begin{center}
\begin{tabular}{crrr}
\toprule
$(b_3,b_5,b_7,b_{11})$ & $R_0+R_1$ & Capacity upper bound & Missed at least\\
\midrule
$(0,0,0,0)$&615&497&118\\
$(0,0,0,1)$&582&482&100\\
$(0,0,1,0)$&562&482&80\\
$(0,0,1,1)$&540&481&59\\
$(0,1,0,0)$&540&484&56\\
$(0,1,0,1)$&524&479&45\\
$(0,1,1,0)$&516&479&37\\
$(0,1,1,1)$&510&488&22\\
\bottomrule
\end{tabular}
\end{center}
Consequently any distinct squarefree system with minimum at least 3 and
all prime factors at most 11 misses at least 22 residues modulo 2310.
In particular, a putative complete cover must use a prime divisor at least
13. This proof uses explicit integer arithmetic, not a solver infeasibility
claim. The analogous minima for supports through $3,5,7$ are respectively
$3,7,18$ missed points in periods $6,30,210$. The bound ceases to exclude
the support through 13. No novelty or priority claim is made for this restriction.

\subsection{Exact joint relocation of an odd/even resource pair}
After constructing a partial system, remove its two classes of moduli
$d$ and $2d$, where $d$ is odd and composite. Let $g_\epsilon(a)$ be the
number of currently uncovered points with parity $\epsilon$ and residue
$a\pmod d$. Replacing the odd class at $a$ gains
$H(a)=g_0(a)+g_1(a)$. An even class at a different odd residue $c$ gains
$G(c)=\max\{g_0(c),g_1(c)\}$. If $c=a$, its additional gain is zero.
Thus the best joint replacement has gain
\[
 \max_{a\ne c}\{H(a)+G(c)\}.
\]
It suffices to consider the two largest entries of $H$ and the two largest
entries of $G$: after forbidding one index, at least one of the top two
remains available. All counts use every integer in the period. Accepting
only strict improvements ensures termination. The parity of each $2p_i$
class is also reoptimized exactly while its odd label stays 1.

\begin{center}
\begin{tabular}{lrrr}
\toprule
Construction & Period & Classes & Missed residues\\
\midrule
Joint solver, primes through 7 & 210 & 14 & 22\\
Joint solver, primes through 11 & 2310 & 30 & 159\\
Joint solver, primes through 13 & 30030 & 62 & 1921\\
After joint pair repairs & 30030 & 62 & 829\\
Extend to 17 and repair all pairs & 510510 & 126 & 8040\\
\bottomrule
\end{tabular}
\end{center}
The solver runs had time limits of 30, 60, and 60 seconds respectively and
returned feasible partial systems, not proofs of optimality. For the
prime-17 experiment, the old classes were first lifted unchanged and
classes modulo $17d$, $d\mid30030$, were added. The old classes were then
allowed to move during repair. The final missed density is
\[
 \frac{8040}{510510}=\frac{268}{17017}=0.015748956925428\ldots.
\]
Every residue of the period was checked directly, and the separate
coordinate-cylinder verifier also rejected the candidate. In particular,
\[
 x\equiv404\pmod{510510}
\]
is missed by all 126 congruences. A small missed density is not a covering
certificate or evidence that the remaining gap is easy to close.

\subsection{Why appending one fresh prime cannot finish this candidate}
For a frozen partial system of period $L$ with residual $R$, a fresh prime
$q$ allows new squarefree moduli $qd$, $d\mid L$. Put
\[
 B(R)=\sum_{d\mid L}\max_{a\bmod d}\#\{x\in R:x\equiv a\pmod d\}.
\]
Over the extended period, there are $q|R|$ residual points. Each class
modulo $qd$ covers at most the corresponding summand of $B(R)$, so a
complete extension requires $q|R|\le B(R)$, regardless of placements.
For the repaired prime-13 candidate, $B(R)=5674$, $|R|=829$; hence no
fresh prime $q\ge17$ can finish it with its old classes fixed.
At $q=17$, the greedy initialization covers 5668 of those residual points,
only six fewer than this total capacity bound, but 8425 points still
remain. Moving the old classes reduces that number to 8040.
For this final candidate, $B(R)/|R|=5093/670<19$. It likewise cannot be
completed by adding just one fresh prime and keeping all its old classes.
This is a restricted extension obstruction, not a theorem excluding a
squarefree minimum-3 cover or joint rearrangements on a larger support.

\subsection{Verification files}
The construction and repair scripts are
\path{code/squarefree_parity_search.py} and
\path{code/squarefree_parity_repair.py}. The strongest partial candidate is
\path{data/squarefree-parity-3-17-repaired.json}, containing literal
congruences and both coverage checks. The counting certificate is
\path{data/squarefree-parity-anchor-audit.json}.
Tests compare 500 joint-pair maximizations with exhaustive alternatives,
check normalization on 100 finite systems, and verify all 31 anchor
capacity cases through prime 13 by direct period enumeration. A squarefree
minimum-2 positive control covers all 210 residues of its period.
These results develop and test a different structural approach; they do
not establish a distinct squarefree covering with minimum modulus 3.

\section{Unrestricted minimum 50: coupled resources and relative arrows}
\label{sec:general50}
The squarefree restriction is dropped in this section. The objective remains
an explicit finite distinct cover, with every modulus at least 50. None of
the following experiments has achieved that objective.

\subsection{An exact two-resource exchange}
In digit coordinates let $C\subseteq F$ be the supports of two resources.
Remove both rows and call the resulting residual $R$. For assignments $a$
on $C$ and $b$ on $F$, write
\[
 H(a)=|R\cap[a]_C|,\qquad G(b)=|R\cap[b]_F|.
\]
The fine cylinder is either contained in the coarse cylinder or disjoint
from it. In the disjoint case the union gain is $H(a)+G(b)$; in the nested
case it is $H(a)$. For any fixed $b$, only one coarse assignment is forbidden
by disjointness. Therefore a globally optimal disjoint pair uses one of the
two largest coarse values. For each such coarse assignment, erase its
cylinder and maximize the fine marginal. Comparing these gains with
$\max H$ gives an exact joint optimum. This argument includes repeated
effective supports caused by a one-state binary tail digit.

The implementation in \path{code/nested_resource_exchange.py} counts the
entire finite kernel using decision diagrams. It accepts only strict
reductions of the uncovered cardinality, and reconstructs the union from
scratch afterwards. Four hundred small cases were compared against all
pairs of placements, including empty residuals, equal supports, and unit
coordinate domains. The existing diagram tests also passed 900 comparisons.

An initial 73-resource target-50 candidate was improved in three exchanges.
It still fails the separate kernel checker and a direct modular witness
check. Its missed progression is
\[
 x\equiv199766327024850602910774556931655400
 \pmod{862639566069660506621770753012536000}.
\]
The search stopped at a time limit; no claim of a local or global optimum
is made. Another prime-stage candidate placed 266 resources and also failed.
These are incomplete constructions, not approximate coverage certificates.

\subsection{Graded supplies of disjoint covering sets}
Fix a target progression $A=a\pmod m$. A supply is a list of sets, each
covering $A$, with pairwise disjoint modulus families. Attach to each set
a grade, its smallest symbolic modulus. Start with the unit set of grade 1.
When introducing a previously unused prime $p$, put $e=v_p(m)$. Keep the
old list, and for $1\leq j\leq e$ make copies constrained to
$x\equiv a\pmod{p^j}$; their grades are multiplied by $p^j$. Partition the
old list into disjoint groups of $p-1$. Each group fills a $p^{e+1}$ arrow
at the prefix $a\pmod{p^e}$, and has grade
\[
 p^{e+1}\min\{g:g\text{ is a grade in the group}\}.
\]
The lists use mutually disjoint $p$-exponent ranges
$0,1,\ldots,e,[e+1,\infty)$. Within the tail range, disjoint groups use
disjoint input spectra. This proves resource distinctness inductively.
Every operation covers $A$; finite closure uses a different auxiliary
prime for each tail occurrence, outside all base primes.

A largest-prime stage $(q,j)$ can now use inputs of grade at least
$\lceil50/q^j\rceil$. At fixed exponent $j$, each input covers one
progression modulo $q^jm$. At an exponent tail starting at $j$, each group
of $q-1$ inputs covers a progression modulo $q^{j-1}m$. Different stages
have disjoint largest-prime/exponent families. If a stage supplies $k$
progressions of the same modulus, their best placement against the current
residual is obtained by choosing the $k$ largest residue counts. Thus a
whole stage is relocated jointly, without sampling.

The experiment used projected period
$L=2^5 3^3 5^2 7=151200$, largest prime 89, and three deterministic
prime-order policies for each shape. It produced 37 guaranteed progressions,
leaving 41308 points outside their union. This is only the gap of the
\emph{guaranteed progression union}; expanded arrow rows can cover additional
points. It is not an uncovered progression for the actual expanded system.

There is no reason to repeat placement searches within this precise menu.
For each stage take the maximum possible guaranteed density over its
available shapes and add these maxima. Exact arithmetic gives
\[
 \sum_{\text{stages}}\max_{\text{searched shapes}}\frac{\min(k,m')}{m'}
 =\frac{4975}{6048}<1.
\]
Here $m'$ is the stage's projected modulus. This excludes a complete union
of guaranteed progressions within the recorded menu, irrespective of their
placements. It does not exclude other prime orders, other stage allocations,
incidental arrow coverage, or minimum 50 itself.

The supply implementation was checked on 48 finite kernels: all 228 input
sets covered their target, had the claimed minimum grade, and had disjoint
exponent resources. A minimum-3 positive control covers every point of its
projected period 360. Duplicate root-exponent stages are rejected.

\subsection{Using coverage inside previously built prime branches}
The next extension retains the individual input sets of the old arrows.
For a prospective target $A$, an old $p$-arrow may already cover $A$ in
$c_p(A)$ of its $p-1$ inputs. A new relative arrow then needs only
\[
 k_p(A)=p-1-c_p(A)
\]
new input sets. For $k_p(A)>0$, the corresponding supply count becomes
\[
 T\longmapsto(v_p(m)+1)T+\left\lfloor T/k_p(A)\right\rfloor.
\]
These are sets covering $A$ relative to the frozen old cover. Their new
moduli become distinct from that cover after multiplication by a fresh
largest root prime. This step must retain and verify the old-coverage
dependencies; counts alone are not a covering certificate.

An initial probe using the explicit old input DAGs examined 342 targets
with modulus at most 720. It found three sufficient input counts for root
97. For example, $33\pmod{720}$ has 96 available relative sets, exactly
the 96 needed inputs. The old guaranteed progression union misses all
210 representatives of this progression modulo 151200. These counts are
saved as a probe, not as a verified replacement or a complete cover.
The subsequent continuation exports and independently audits these schemas,
as described below, and uses branch information during the main construction.

\paragraph{Reproducible records.}
The algorithms are \path{code/graded_arrow_blocks.py} and
\path{code/borrowed_arrow_probe.py}. Results, exact menu capacities, and
the pending relative-input witnesses are in
\path{data/construction-50-graded-blocks.json} and
\path{data/construction-50-borrowed-probe-97.json}. The active research
objective is still open; none of these files claims a minimum-50 cover.

\subsection{Verified relative schemas and a stronger early core}
The count witness for $33\pmod{720}$ has now been converted into an explicit
relative-arrow DAG and independently checked. It uses root 97, has 96
inputs, and contains 29 references to already covered inputs of earlier
prime arrows. Its finite completion covers $33\pmod{720}$ relative to the
old system, reducing the gap of the guaranteed progression union from
41308 to 41098 modulo 151200. This is a verified partial replacement, not
a covering of all integers.

For clarity, the dependency induction is as follows. Let $B$ be the union
of all earlier largest-prime stages. A new input set must cover $A\setminus B$.
At an internal prime $p$, any omitted digit has a recorded earlier $p$-stage
input that covers $A$, possibly relative to still earlier stages. Thus
$A\setminus B$ has no points in that digit. Each remaining digit receives
a new input set covering $A\setminus B$. Fixed copies and disjoint tail
groups preserve disjointness of the \emph{new} modulus families; borrowed
classes remain in $B$ and are not copied into the new system. Multiplication
by the new largest prime separates the new families from all old families.
References always go to earlier largest-prime stages, so the induction is
acyclic. Unique auxiliary primes close all own tail occurrences finitely.

Two implementations check whether an earlier input covers a requested
progression. The construction uses a first-nonzero-digit calculation.
The independent audit enumerates compatible CRT intersections:
\[
 a\equiv r+d p^{k-1}\pmod{\gcd(m,p^k)},\qquad
 E\leq k\leq\max(E,v_p(m)+1).
\]
Here $r$ is the arrow prefix and $E$ its first exponent. One level beyond
$v_p(m)$ represents every later possible first-nonzero level. The audit
also reconstructs every recorded supply operation, verifies the root
inputs' grades, and checks that no stage's exponent family is reused.

An integrated construction rebuilds the stages above prime 7 with these
dependencies. When there are spare inputs, it prefers smaller admissible
grades, preserving broader coverage for later stages. Through root 113,
with target moduli at most 720, it obtains 22 new root stages and 708
checked borrowed inputs. The guaranteed union leaves 35149 residues.
Allowing roots through 229 and target moduli at most 2520 produces 36
stages, the last at root 173, with 2119 checked borrowed inputs and a gap
of 33689 modulo 151200. Both are verified partial schemas. In particular,
the first projected gap is not claimed to miss all their expanded rows.

The new dependency tests exhaust all 384 states of a small relative-arrow
example. They verify 405 containment comparisons, 7008 borrowed-coverage
implications, and disjointness of all 12 own exponent families; the combined
small example covers every state. These tests supplement the induction
and the independent audit, rather than replacing them.

The early core is now being improved separately. An exact array calculation
places 283 resources on the prime set $\{2,3,5,7\}$ with depths
$(6,4,3,2)$ and resource bases at most 50000. It checks all 7257600 kernel
states, leaves 2998392 uncovered, and retains a separately verified modular
counterexample. The uncovered fraction is
\[
 \frac{2998392}{7257600}=\frac{124933}{302400}.
\]
This uses partial prime-tail coverage that the progression-only core loses.
The region-based extension below integrates this core into the
relative-arrow construction, with a separate dependency audit.

The verified records are \path{data/construction-50-relative-97-certified.json},
\path{data/construction-50-relative-113.json}, and
\path{data/construction-50-relative-229.json}; the core is in
\path{data/construction-50-exact-core.json}. Implementations are
\path{code/relative_arrow_search.py} and \path{code/exact_core_array.py}.

\subsection{Exact residual regions and a complete minimum-20 control}
The implementation \path{code/core_region_arrows.py} now allows targets
which fix a core tail digit, and computes their coordinate supports from
the entire residual array. At each stage it may use either the frozen core
residual or the residual left after all previously certified target regions.
The latter option requires chronological verification: a region may be
removed only after its own relative-cover implication has been checked.
When checking a child inside a prime branch, the stronger option conditions
the actual residual on that branch before taking the other coordinate
supports. This preserves correlations that independent marginal supports
would lose.

The induction remains a relative covering argument. Let $B$ contain the
core and all earlier stages. Every new root input must cover the target
outside $B$. Omitted core digits are absent from the exhaustively computed
residual, and omitted noncore digits have recorded earlier input witnesses.
New own exponent families remain disjoint. Each root uses a new largest
structural prime; finite closure uses distinct fresh auxiliary primes.
The checker independently tests branch membership by congruences, rebuilds
the supply operations, and verifies all grades and backward dependencies.

A minimum-20 control with core depths $(5,3,2,2)$ starts with 15995 missing
states out of 241920. Sixteen root stages, from 11 through 71, close this
gap completely. Its audit checks 370 borrowed inputs and returns
\texttt{VERIFIED\_COMPLETE\_REGION\_SCHEMA}. The record is
\path{data/construction-region-control-20.json}. This is a finite-completion
schema, not an enumerated list of its potentially enormous finite expansion.
It is a correctness control, not a new record or completion of the active
minimum-50 objective.

For minimum 50, concentrating the early core on the even class gives a
stronger combined result than the core that minimizes its own gap alone.
The early core leaves 3419270 states, but the subsequent 26 root stages
through 113, with 1529 checked borrowed inputs, reduce the guaranteed gap to
\[
  \frac{812946}{7257600}\simeq 0.112013.
\]
This is the fraction of finite kernel states outside the guaranteed region
union; it is neither an integer density nor a demonstrated counterexample
to every expanded congruence. The result is only
\texttt{VERIFIED\_PARTIAL\_REGION\_SCHEMA}, saved in
\path{data/construction-50-focus-2-regions-113.json}.
The full-state tests check 6564 marginal comparisons and 5760 combined
core/noncore states; the branch-sensitive control checks 7680 relative
coverage implications and rejects deliberately false certificates.

\subsection{Union inputs and simultaneous resource packing}
The next structural extension lets a single input be a union of several
prime-power cylinders. For example,
\[
  (1\bmod 3)\ \cup\ (3\bmod 9)
\]
can use one exponent-1 resource and one exponent-2 resource. A rule which
only groups equal-tail inputs cannot represent this allocation directly.
The program \path{code/prime_forest_supply.py} enumerates antichain covers
of a prescribed one-prime state set. Each cover has a count vector recording
how many old inputs it needs at each exponent level. Constructive packing
allocates each pair $(\hbox{exponent},\hbox{old input})$ at most once.
Old inputs at exponent zero are retained. Because old input families are
disjoint, this proves disjointness of the resulting new modulus families.
The packing heuristic is not asserted to be optimal.

The region checker verifies each recorded forest directly: every needed
coordinate state is covered, every branch uses a valid old child, no
exponent/child pair is repeated, and the minimum grade is correct. Its
small exhaustive control checks 8640 relative implications. The primitive
packing tests check 324 complete antichain covers in 64 cases. Target
regions may also unite a residue class modulo $p$ with a different class
modulo $p^2$; their counts are computed exactly.

These first forest runs have not improved the best minimum-50 result.
The bounded run through root 53 leaves 1058058 guaranteed missing states;
the run permitting union targets stops after root 29 with 1492358 missing
states. Neither is a complete cover. Their records are
\path{data/construction-50-forest-regions-113.json} and
\path{data/construction-50-union-regions-113.json}; the filename's 113 is
the allowed root ceiling, not the last root reached before the time limit.

A further implementation, \path{code/joint_root_repair.py}, chooses all
$q-1$ inputs simultaneously from a shared pool of fixed-exponent resources.
Each input may combine cylinders in several primes. Each resource may be
used at most once across all inputs, while every projected residual point
must be covered in every input. This projection is exhaustive: a point is
included whenever its whole fiber contains any core-uncovered target state.
Thus a successful assignment proves the relative target implication.
The checker independently reconstructs the projection, checks every input,
and verifies global resource uniqueness and the minimum modulus.

For the target
\[
 (1\bmod4)\cap\bigl((1\bmod3)\cup(3\bmod9)\bigr)
\]
in the minimum-50 Owens-inspired core, the first root-11 test has 784
projected points, 60 resource types, and 4315 candidate classes. The bounded
solver run returns \texttt{UNKNOWN}; it provides neither a replacement nor
an impossibility proof. A separate positive control passes direct checks
over a whole integer period and rejects duplicated resources and uncovered
inputs. The record is \path{data/construction-50-joint-root-11.json}.

A subsequent exact capacity audit resolves this restricted attempt. For
nonnegative integer weights $w_x$ on the projected target, every disjoint
collection of $t$ complete inputs must satisfy
\[
 t\sum_x w_x\ \leq\ \sum_{\mathbf e}
       \max_{\mathbf r}\sum_{x\in C(\mathbf e,\mathbf r)}w_x.
\]
A numerical linear program proposes weights; an independent Python integer
histogram calculation then verifies the displayed inequality's right side.
For this 60-resource menu the saved weights give
\[
 t\leq\frac{6351876232}{1000000179}<10.
\]
Thus this particular fixed-exponent menu cannot fill the ten root-11 inputs,
regardless of how its classes are combined. The certificate is
\path{data/construction-50-joint-root-11-capacity.json}. It does not exclude
tail families, deeper exponent resources, different targets, or a
minimum-50 cover. The following extension implements tail-aware joint projection.

\subsection{Auditing the early Owens-inspired seed}
The explicit seed in \path{code/owens_seed_core.py} includes chosen
placements inspired by the early $2,3,5,7$ stages. It does not purport to
transcribe or verify the full thesis. An independent array audit checks all
7257600 core states. At cutoff 42 it verifies the three local claims:
$8\bmod16$ is filled; $4\bmod8$ is filled outside $1\bmod5$; and
$16\bmod32$ is filled outside $4\bmod5$.
Raising this seed's cutoff to 50 removes its three resources with bases
42, 45, and 48. The respective local gaps above then become 17160, 34752,
and 3500 states. These are overlapping local diagnostics, not a partition
of the whole residual. The saved audit is
\path{data/owens-seed-branch-audit.json}. The active objective remains open:
none of the minimum-50 records is a verified complete covering system.

\subsection{Tail-aware joint inputs and a verified local 42-branch repair}
The joint-input model now permits a terminal tail resource in each core
coordinate. To project to a lower tail depth $E$, it keeps the prefix modulo
$p^{E-1}$ and the first nonzero digit after that prefix. This is determined
by every state of the finer core kernel. The projected target contains a
state if any fine state in its fiber remains uncovered, so covering the
projected target is sufficient for covering the entire original residual.
The independent test computes the same projection by congruence enumeration.

This model finds a concrete root-11 repair of
\[
 A=10\pmod{42}
  =\{x:x\equiv0\pmod2,\ x\equiv1\pmod3,\ x\equiv3\pmod7\}
\]
relative to the saved minimum-50 early core inspired by Owens.
It uses local tail depths $(2,2,1,2)$ and 40 distinct exponent resources
distributed among ten inputs. Their minimum grades, in input order, are
\[
 7,9,5,14,21,15,30,45,6,10.
\]
Consequently every new modulus is at least $11\cdot5=55$.
All 40 projected target states are covered in every input; these account
for 102240 uncovered fine core states in $A$. The independently checked
record is \path{data/construction-50-joint-tail-repair-42.json}.
Unique auxiliary-prime closure makes these symbolic tails finite while
preserving their resource distinctions. This repairs a local branch, not
Owens's entire covering: it consumes prime-11 resources that his later
construction needs elsewhere.

The local repair can now initialize the relative-region construction.
Mixed-prime input unions are checked directly on the actual conditioned
core residual, allowing later stages to borrow their coverage.
The independent small test checks 144 projection states and 6888 borrowing
implications, and its integrated control completes a cover.
For minimum 50, the bounded continuation reaches root 61 with 14 stages,
including the joint root, and 160 checked borrowed inputs. It still leaves
1487500 guaranteed missing core states. The record is
\path{data/construction-50-repair42-continuation.json}.

Increasing the joint root-11 menu for the old Nielsen target does not by
itself resolve that target. Exact integer-weight certificates give
\[
 \frac{8975856116}{1000000087}<10
 \quad\hbox{at depths }(3,3,2,1),\qquad
 \frac{9125393259}{999999621}<10
 \quad\hbox{at depths }(3,4,2,1).
\]
These bound only the respective saved menus; they do not exclude deeper
resources, other early cores, or other targets.

Another test asks one root-11 stage to repair the union of the three early
classes with bases 42, 45, and 48. At depths $(5,3,2,2)$ its target has
12442 projected states and 212 eligible resources. The explicit class
model has 174560 candidates and returns \texttt{UNKNOWN} in its bounded
run. A second implementation chooses residue labels implicitly and adds
exact uncovered-state witnesses after checking every projected state.
It reaches 24 witness constraints without finding a verified replacement.
No positive claim is made from intermediate witness solutions. See
\path{data/construction-50-joint-three-hole-root-11.json} and
\path{data/construction-50-joint-three-hole-refinement.json}.

\subsection{A fully audited transcription of the published prime-11 stage}
To compare with the actual published mechanism, the ten inputs from
Nielsen's Section 4.5 have been written explicitly, with Owens's stated
prime-5 permutation. The transcription uses translated tail centers:
the relevant 27-tail is centered at 21, whereas the full 81-tail in the
eighth input is centered at 3. These are distinct branches; confusing them
leaves an exactly detectable hole.

The corrected interpretation in \path{code/nielsen_prime11.py} passes an
exhaustive relative audit against the explicit cutoff-42 early seed.
Its 417 expanded exponent resources are all distinct, and each of its ten
inputs covers all 628992 target states left by that seed. The record
\path{data/nielsen-prime11-transcription.json} is a verified local stage,
not a verification of the full thesis.

The cutoff-50 audit then identifies the precise damage. Input 1 contains
the inadmissible grade-4 family, whose first root modulus is $11\cdot4=44$.
Inputs 6 and 7 each miss 28000 target states after the early modulus-45
class is removed; inputs 9 and 10 each miss 5600. Inputs 2--5 and 8 retain
their relative coverage. This diagnostic is saved in
\path{data/nielsen-prime11-minimum50-diagnostic.json}.
It motivates reallocating resources among these particular inputs and
retaining valid higher prime-power parts of a family when only its first
modulus falls below 50. Neither this diagnosis nor the local 42-branch
repair is a complete minimum-50 covering system.

\subsection{Partial borrowing at prime 13 and exact input transport}
The prime-13 step is now explicit as well. Ten inputs are translated copies
of the prime-11 pattern, with Owens's prime-5 permutation. Two more inputs
use an internal 11-tail: the new rows on each digit are combined with the
\emph{own} rows already present in that digit of the earlier prime-11 stage.
Neither contribution has to cover the entire core target by itself.
Borrowed rows are never copied into the new modulus family.

This reproduces the published doubling mechanism. In the first new copy,
the binary-4 and binary-8-tail factors are transported from $1\bmod4$ to
$3\bmod4$. In the second, those factors are replaced by 1 and the odd
class modulo 2. The first copy uses binary exponents at least 2 and the
second uses exponents 0 and 1, so their new families are disjoint.
Checking the union of old and new coverage separately on every 11-digit
verifies all 837 expanded resource families; the smallest new modulus is
52. This is the local record \path{data/nielsen-prime13-transcription.json}.

A compact equivalent representation uses depths $(3,4,2,1)$ for the input
bodies. It has 117 resources in the prime-11 stage and 237 in the prime-13
stage. Its independent audit checks 68190 prime-13 state conditions.
The construction checker now authenticates the older input references,
checks the new resource distinctions, and evaluates these mixed-prime
unions on the actual residual. The cached union evaluator agrees with
direct evaluation in 600 test cases; deliberately repeated resources and
incorrect transported prefixes are rejected. Files are
\path{data/nielsen-prime11-compact.json},
\path{data/nielsen-prime13-compact.json}, and
\path{data/partial-borrow-transport-tests.json}.

The cutoff-42 continuation from these two audited stages remains partial:
through root 113 it leaves 521808 guaranteed missing states. This checks
the stronger mechanism but does not verify the remainder of Owens's thesis.

\subsection{Relocating modulus 90 restores the odd input targets at minimum 50}
Keeping the early core fixed led to bounded but unsuccessful searches for
exchanges among the damaged prime-11 inputs. A simpler structural trade
works: replace the existing core class
\[
 78\pmod{90}\quad\hbox{by}\quad3\pmod{90}.
\]
The removed modulus-45 class was $3\pmod{45}$. Its odd part is exactly
$3\pmod{90}$, whereas the old $78\pmod{90}$ class is even. Therefore this
trade restores all modulus-45 coverage needed on
\[
 T=\{x:x\text{ is odd}\}\cap
       \bigl((1\bmod3)\cup(3\bmod9)\bigr).
\]
The deleted core rows with bases 42 and 48 do not meet $T$. Consequently
the cutoff-50 relocated core has \emph{exactly the same residual on $T$}
as the cutoff-42 core. An exhaustive comparison of all 1612800 kernel
states in $T$ verifies this equality; both residuals contain 1257984 states.
The lost even class remains an explicit repair obligation.

This gives nine complete admissible prime-11 inputs without changing their
published resource pattern. The nine-input certificate uses 116 distinct
resources, and its audit checks 20457 input/state pairs. It is intentionally
labelled \texttt{VERIFIED\_INPUT\_PACKING}: nine inputs do not fill all ten
digits of an 11-tail. The default root verifier rejects that incomplete
packing as a complete root. Records are
\path{data/construction-50-relocated90-core.json},
\path{data/construction-50-relocated90-odd-agreement.json}, and
\path{data/construction-50-relocated90-nine-inputs.json}.

The inadmissible first input need not be discarded wholesale. Keep its
binary class $1\bmod4$ and all root-11 powers starting with $11^2$; delete
only the first class, of modulus 44. Together with the other nine inputs,
this gives 233 disjoint exponent resources of minimum modulus at least 50.
On the prime-11 target, 227300 of 250030 projected core/root states are
covered. The remaining 22730 states are precisely the part with
$x\equiv1\pmod{11}$; since the target already fixes $x\equiv1\pmod4$,
the exception is contained in $1\pmod{44}$. This is the relative statement
in \path{data/construction-50-relocated90-prime11-truncated.json}, not a
claim about every integer outside that class.

The prime-13 target lies in $3\bmod4$, so the missing modulus-44 class
contributes nothing there. Its complete relative stage remains valid with
the relocated minimum-50 core and the nine admissible earlier inputs.
The checker explicitly uses an empty first earlier input when authenticating
this proof. The resulting minimum-50 record is
\path{data/construction-50-relocated90-prime13.json}.

Finally, the supply search now interleaves usable old primes with core
primes according to their effective input multiplier
\[
 f+1+\frac1k,
\]
where $f$ is the number of available fixed-prefix copies and $k$ is the
number of tail digits still requiring new inputs. This permits a cheap
partially filled 11-tail to be used before primes 5 or 7, rather than
forcing all core primes to precede it. Combined with the relocated core,
the root-17 stage closes the outstanding part of $T$ with 23 available
admissible inputs. The bounded continuation through root 59 has 13 root
stages and 160 checked borrowed inputs, but still leaves 1314132 guaranteed
missing states. It is saved in
\path{data/construction-50-relocated90-continuation.json}.
No complete minimum-50 covering system has yet been obtained. The remaining
work includes the even holes created or exposed by the core changes and
more effective use of partial coverage in subsequent stages.

\begin{thebibliography}{9}
\bibitem{daltontrifonov} J. Dalton and O. Trifonov,
\emph{Extreme Covering Systems}, Journal of Integer Sequences
\textbf{25} (2022), Article 22.9.1, especially Problem 4.
\url{https://cs.uwaterloo.ca/journals/JIS/VOL25/Trifonov/trif3.pdf}.
\bibitem{cft} M. Cummings, M. Filaseta, and O. Trifonov,
\emph{An upper bound for the minimum modulus in a covering system with
squarefree moduli}, arXiv:2211.08548 (2022), Theorem 1.1.
\url{https://arxiv.org/abs/2211.08548}.
\bibitem{nielsen} P. P. Nielsen,
\emph{A covering system whose smallest modulus is 40},
J. Number Theory \textbf{129} (2009), 640--666.
\href{https://doi.org/10.1016/j.jnt.2008.09.016}{doi:10.1016/j.jnt.2008.09.016}.
\bibitem{owens} T. Owens,
\emph{A Covering System with Minimum Modulus 42},
MS thesis, Brigham Young University, 2014.
\url{https://scholarsarchive.byu.edu/etd/4329/}.
\bibitem{agrawal} R. Agrawal, P. Bhatia, K. Gupta, P. Lamb, A. Lott,
A. Rice, and C. R. Ward,
\emph{Computations and observations on congruence covering systems},
Integers \textbf{24A} (2024), A1.
\url{https://math.colgate.edu/~integers/a1Proc23/a1Proc23.pdf}.
\end{thebibliography}
\end{document}
